% SOURCE_AUTHORITY: sga5_fr_workpass.tex, lines 8288--8338 (LF-normalized unit).
% SOURCE_SHA256: 791F4EFFC5E02832D5D77ED03518C8156D6F07E4C8238B03545DB93D883FBB28
Se llamará categoría de los \(\mathbb Z_\ell\)-haces constantes torcidos constructibles a la subcategoría plena de \(\mathbb Z_\ell\text{-}\fc(X)\) generada por los \(\mathbb Z_\ell\)-haces constantes torcidos constructibles.

\begin{statement}{Ejemplo 1.2.2}
El \(\mathbb Z_\ell\)-haz \(\mathbb Z_\ell(1)\).
\end{statement}

Supongamos que el entero \(\ell\) es primo con las características residuales de \(X\). Para todo entero \(n\), la elevación a la potencia \(\ell\) define un morfismo de haces localmente constantes:
\[
\mu_{\ell^{n+1}}\longrightarrow \mu_{\ell^n}.
\]
Esto define de manera evidente un sistema proyectivo:
\[
\mathbb Z_\ell(1)=(\mu_{\ell^{n+1}})_{n\in\mathbb N}.
\]
Reduciéndose, para cada entero \(n\), al caso en que \(\mu_{\ell^n}\) y \(\mu_{\ell^{n+1}}\) son constantes, se comprueba que el sistema proyectivo \(\mathbb Z_\ell(1)\) es \(\ell\)-ádico.

Cuando \(X=\Spec(K)\), con \(K\) un cuerpo de tipo finito sobre el cuerpo primo, el \(\mathbb Z_\ell\)-haz \(\mathbb Z_\ell(1)\) es un ejemplo de haz \(\ell\)-ádico constante torcido constructible (1.2.1) para el cual no existe ninguna extensión finita étale \(K'\) de \(K\) tal que las imágenes inversas de los \(F_n\) sobre \(X'=\Spec(K')\) sean constantes. En efecto, \(K'\) contendría, para todo entero \(n\), una raíz primitiva \(\ell^n\)-ésima de la unidad, lo que no es posible para una extensión de tipo finito del cuerpo primo.

\begin{statement}{Recordatorios 1.2.3}
(SGA 1 V 7). Se sabe que la categoría de los haces abelianos de torsión constructibles localmente constantes sobre un preesquema \(T\) es equivalente a la categoría de los recubrimientos étale de \(T\) dotados de una estructura de \(X\)-grupo abeliano.
\end{statement}

Sea ahora \(X\) un preesquema localmente noetheriano conexo. Denotemos por \(R_X\) la categoría de los recubrimientos étale de \(X\) y por \(\Ensf\) la de los conjuntos finitos. Si \(a\) es un punto geométrico de \(X\), sea \(F_a\) el funtor
\[
F_a:R_X\longrightarrow \Ensf,
\]
que a todo recubrimiento étale de \(X\) asocia su fibra geométrica en el punto \(a\). Como el grupo fundamental
\[
\pi_1(X)=\pi_1(X,a)
\]
de \(X\) en el punto \(a\) es, por definición, el grupo de los automorfismos de \(F_a\), convenientemente topologizado, el funtor \(F_a\) puede interpretarse como un funtor
\[
R_X\longrightarrow \Ensf(\pi_1)
\]
de la categoría \(R_X\) en la categoría de los conjuntos finitos dotados de una acción continua de \(\pi_1=\pi_1(X,a)\).

Recordemos la proposición siguiente:

\begin{statement}{Proposición 1.2.3.1}
El funtor
\[
F_a:R_X\longrightarrow \Ensf(\pi_1)
\]
es una equivalencia de categorías. Conmuta con los límites inductivos y con los límites proyectivos finitos.
\end{statement}

Sean ahora \(Lc(X)\) y \(\Abf(\pi_1)\), respectivamente, la categoría de los haces abelianos de torsión localmente constantes constructibles y la de los grupos abelianos finitos dotados de una acción continua del grupo fundamental. El funtor \(F_a\) induce un funtor exacto, denotado también por \(F_a\),
\[
F_a:Lc(X)\longrightarrow \Abf(\pi_1),
\]
que es una equivalencia de categorías.
